\chapter{Fonctions affines et comparaisons}
\section*{Exercice 13.1 -- Reconstruire}
La pente vaut $(14-5)/(4-1)=3$. Puis $p=5-3\times1=2$, donc $f(x)=3x+2$. Les deux données sont vérifiées.

\section*{Exercice 13.2 -- Signe et variation}
La pente $-4$ est négative, donc $f$ est strictement décroissante sur $\R$. Son zéro est $3$ ; elle est positive pour $x<3$, nulle en $3$, négative pour $x>3$. Les deux études ne doivent pas être confondues.

\section*{Exercice 13.3 -- Comparer des tarifs}
$A-B=20-4x$. Les tarifs sont égaux pour $x=5$ ; $B$ est moins cher pour $0\leq x<5$, $A$ pour $x>5$. À $5$, chacun coûte $40$.

\section*{Exercice 13.4 -- Une preuve}
Pour $a<b$, $f(b)-f(a)=m(b-a)$. Les deux facteurs étant positifs, la différence l'est aussi, donc $f(a)<f(b)$. Les entrées ont été choisies quelconques : la preuve porte sur tout le domaine.

\section*{Exercice 13.5 -- Interpolation}
La pente est $(30-18)/(8-2)=2$, d'où $f(t)=18+2(t-2)$. À $t=5$, elle vaut $24$. Le résultat suppose une variation affine entre les relevés ; les deux mesures seules ne l'imposent pas.

\section*{Exercice 13.6 -- Modèle physique}
$V(t)=75-2{,}5t$ sur $[0;30]$. L'équation $75-2{,}5t=20$ donne $t=22$ minutes, valeur admissible. Le modèle s'arrête à la vidange, $t=30$.
